\section*{Chapitre \ref{ch:g5:water-shapes-the-land} --- L'eau façonne les paysages}

\begin{solution}{exo:g5:water-shapes-the-land:1}
(a) \omterm{def:g5:water-shapes-the-land:weathering}{Altération}~: la \omterm{def:g3:stones-pebbles-sand:rock}{roche} est brisée sur place. (b) \omterm{def:g5:water-shapes-the-land:erosion}{Érosion}~: le \omterm{def:g3:stones-pebbles-sand:sand}{sable} est
arraché et emporté. (c) \omterm{def:g5:water-shapes-the-land:deposition}{Dépôt}~: la \omterm{def:g3:stones-pebbles-sand:sand}{boue} se dépose. (d) \omterm{def:g5:water-shapes-the-land:weathering}{Altération}~: la
racine fend la pierre sur place.
\end{solution}

\begin{solution}{exo:g5:water-shapes-the-land:2}
La \omterm{def:g5:water-shapes-the-land:source-mouth}{source} est l'endroit où la \omterm{def:g1:landscapes-around-us:water}{rivière} commence (une \omterm{def:g5:water-shapes-the-land:source-mouth}{source} d'eau, un
glacier qui fond, une terre haute et humide). L'\omterm{def:g5:water-shapes-the-land:source-mouth}{embouchure} est l'endroit où
elle se termine, dans la \omterm{def:g2:water-on-earth:ocean}{mer}, dans un \omterm{def:g1:landscapes-around-us:water}{lac} ou dans une \omterm{def:g1:landscapes-around-us:water}{rivière} plus grande.
Un \omterm{def:g5:water-shapes-the-land:source-mouth}{affluent} est une \omterm{def:g1:landscapes-around-us:water}{rivière} plus petite qui se jette dans une plus grande.
\end{solution}

\begin{solution}{exo:g5:water-shapes-the-land:3}
La \omterm{def:g5:water-shapes-the-land:source-mouth}{source}~; la \omterm{def:g1:landscapes-around-us:landform}{vallée} étroite en V~; les \omterm{def:g5:water-shapes-the-land:meander}{méandres}~; le \omterm{def:g5:water-shapes-the-land:delta}{delta}.
\end{solution}

\begin{solution}{exo:g5:water-shapes-the-land:4}
(a) Faux~: le gel, l'eau, les racines, la chaleur et le froid brisent
lentement même la \omterm{def:g3:stones-pebbles-sand:rock}{roche} dure (\omterm{def:g5:water-shapes-the-land:weathering}{altération}). (b) Vrai~: plus l'eau est
rapide, plus les morceaux qu'elle peut déplacer sont gros. (c) Faux~: les
morceaux les plus lourds se déposent en premier~; la \omterm{def:g3:stones-pebbles-sand:sand}{boue} la plus fine se
dépose en dernier, là où l'eau est presque immobile, dans un \omterm{def:g1:landscapes-around-us:water}{lac} ou dans la
\omterm{def:g2:water-on-earth:ocean}{mer}.
\end{solution}

\begin{solution}{exo:g5:water-shapes-the-land:5}
A est à l'extérieur de la courbe~: l'eau y est rapide et profonde, et elle
use la berge (\omterm{def:g5:water-shapes-the-land:erosion}{érosion}). B est à l'intérieur de la courbe~: l'eau y est
lente, et elle dépose du \omterm{def:g3:stones-pebbles-sand:sand}{sable} (\omterm{def:g5:water-shapes-the-land:deposition}{dépôt}).
\end{solution}

\begin{solution}{exo:g5:water-shapes-the-land:6}
En \omterm{def:g1:landscapes-around-us:landform}{montagne}, la \omterm{def:g1:landscapes-around-us:water}{rivière} est pentue et rapide~: elle peut rouler de gros
cailloux, et elle emporte les petits grains. Près de la \omterm{def:g2:water-on-earth:ocean}{mer}, la \omterm{def:g1:landscapes-around-us:water}{rivière} est
lente~: elle ne peut plus déplacer de gros cailloux, seulement du \omterm{def:g3:stones-pebbles-sand:sand}{sable} et
de la \omterm{def:g3:stones-pebbles-sand:sand}{boue}, et elle les dépose là. En chemin, de plus, les cailloux ont été
brisés et usés en grains plus petits.
\end{solution}

\begin{solution}{exo:g5:water-shapes-the-land:7}
(a) Cours supérieur~: $600 \div 6 = 100$ kilomètres. Cours moyen~:
$600 \div 2 = 300$ kilomètres. (b) Cours inférieur~: $600 - 100 - 300 = 200$
kilomètres. Comme $3 \times 200 = 600$, il représente $\frac{1}{3}$ de la
\omterm{def:g1:landscapes-around-us:water}{rivière}. (Vérification~: $\frac{1}{6} + \frac{1}{2} + \frac{1}{3} = \frac{1}{6} +
\frac{3}{6} + \frac{2}{6} = \frac{6}{6}$.)
\end{solution}

\begin{solution}{exo:g5:water-shapes-the-land:8}
Les taches claires sont du \omterm{def:g3:stones-pebbles-sand:sand}{sable} et des \omterm{def:g3:stones-pebbles-sand:pebble}{graviers} déposés à l'intérieur des
courbes, là où l'eau est lente (\omterm{def:g5:water-shapes-the-land:deposition}{dépôt}). Les \omterm{def:g1:landscapes-around-us:water}{lacs} sombres en croissant sont
d'anciennes boucles de la \omterm{def:g1:landscapes-around-us:water}{rivière}~: deux boucles se sont tellement
rapprochées que la \omterm{def:g1:landscapes-around-us:water}{rivière} a percé le pédoncule de terre qui les séparait,
a pris le raccourci et a laissé l'ancienne boucle en arrière, sous la forme
d'un \omterm{def:g1:landscapes-around-us:water}{lac}.
\end{solution}

\begin{solution}{exo:g5:water-shapes-the-land:9}
Les prés plats qui bordent la \omterm{def:g1:landscapes-around-us:water}{rivière} sont sa \omterm{def:g5:water-shapes-the-land:flood}{plaine d'inondation}~: la
terre que la \omterm{def:g1:landscapes-around-us:water}{rivière} recouvre quand elle est en \omterm{def:g5:water-shapes-the-land:flood}{crue}. Ils sont plats
justement parce que la \omterm{def:g1:landscapes-around-us:water}{rivière} les a construits avec sa propre \omterm{def:g3:stones-pebbles-sand:sand}{boue}. Un
jour, la \omterm{def:g1:landscapes-around-us:water}{rivière} sera de nouveau en \omterm{def:g5:water-shapes-the-land:flood}{crue} et remplira l'école d'eau, et une
eau de \omterm{def:g5:water-shapes-the-land:flood}{crue} rapide est dangereuse. Il faut construire l'école sur une
hauteur, et garder les prés comme champs, comme prés ou comme parc.
\end{solution}

\begin{solution}{exo:g5:water-shapes-the-land:10}
(a) Une année compte 12 mois~: $12$ milliards $\div\, 12 = 1$ milliard de
tonnes par mois en moyenne. (b) Un milliard, c'est $1\,000\,000\,000$~;
divisé par 10, cela fait $100\,000\,000$~: cent millions de chargements de
camion.
\end{solution}

\begin{solution}{exo:g5:water-shapes-the-land:11}
Une histoire possible. Le grain fait partie d'une \omterm{def:g3:stones-pebbles-sand:rock}{roche}, haut sur une
\omterm{def:g1:landscapes-around-us:landform}{montagne}. Nuit après nuit, le gel ouvre une fissure autour de lui, jusqu'à
ce qu'un morceau de \omterm{def:g3:stones-pebbles-sand:rock}{roche} se détache~: c'est l'\omterm{def:g5:water-shapes-the-land:weathering}{altération}. Le morceau tombe
dans le ruisseau en contrebas. Dans le cours supérieur, l'eau est rapide et
la pente raide~: elle roule le morceau (\omterm{def:g5:water-shapes-the-land:erosion}{érosion}), et, en se cognant contre
d'autres cailloux, il se brise en \omterm{def:g3:stones-pebbles-sand:pebble}{galets}, puis en grains de \omterm{def:g3:stones-pebbles-sand:sand}{sable}. Dans le
cours moyen, la \omterm{def:g1:landscapes-around-us:water}{rivière} ralentit~; le grain se repose un moment sur la
plage de \omterm{def:g3:stones-pebbles-sand:sand}{sable} à l'intérieur d'une courbe de \omterm{def:g5:water-shapes-the-land:meander}{méandre}, jusqu'à ce qu'une
\omterm{def:g5:water-shapes-the-land:flood}{crue} le reprenne. Dans le cours inférieur, la \omterm{def:g1:landscapes-around-us:water}{rivière} est lente et large~;
le grain voyage avec l'eau jusqu'à l'\omterm{def:g5:water-shapes-the-land:source-mouth}{embouchure}. Là, l'eau s'arrête
presque, et le grain se dépose avec la \omterm{def:g3:stones-pebbles-sand:sand}{boue} et le \omterm{def:g3:stones-pebbles-sand:sand}{sable} du \omterm{def:g5:water-shapes-the-land:delta}{delta}~: c'est le
\omterm{def:g5:water-shapes-the-land:deposition}{dépôt}.
\end{solution}
